\documentclass[10pt,a4paper]{amsart}
\usepackage{amsmath,amssymb,amsthm,graphicx}
% Retained preamble for the Stacks Project chapter morphisms.tex (Morphisms of Schemes).
% The chapter file pinned-original-source.tex carries only \input{preamble} on its first line, so no
% preamble text lives inside the chapter itself.  The Stacks Project uses one single shared preamble for
% every chapter; the snapshot stages it verbatim in sources/stacks-chow/preamble.tex.  The block below is
% copied mechanically, line for line, from that file's lines 14-18 (the Xy-pic setup) and lines 154-217
% (the native theorem declarations together with the mathematical macros).  Nothing was added, removed or
% redefined.  The wrapper class and the amsmath/amssymb/amsthm/graphicx base packages are supplied by the
% assembler, so the class/\externaldocument part of the shared preamble is not repeated here.
% For commutative diagrams we use Xy-pic
\usepackage[all]{xy}

% We use 2cell for 2-commutative diagrams.
\xyoption{2cell}
% Theorem environments.
%
\theoremstyle{plain}
\newtheorem{theorem}[subsection]{Theorem}
\newtheorem{proposition}[subsection]{Proposition}
\newtheorem{lemma}[subsection]{Lemma}

\theoremstyle{definition}
\newtheorem{definition}[subsection]{Definition}
\newtheorem{example}[subsection]{Example}
\newtheorem{exercise}[subsection]{Exercise}
\newtheorem{situation}[subsection]{Situation}

\theoremstyle{remark}
\newtheorem{remark}[subsection]{Remark}
\newtheorem{remarks}[subsection]{Remarks}

\numberwithin{equation}{subsection}

% Macros
%
\def\lim{\mathop{\mathrm{lim}}\nolimits}
\def\colim{\mathop{\mathrm{colim}}\nolimits}
\def\Spec{\mathop{\mathrm{Spec}}}
\def\Hom{\mathop{\mathrm{Hom}}\nolimits}
\def\Ext{\mathop{\mathrm{Ext}}\nolimits}
\def\SheafHom{\mathop{\mathcal{H}\!\mathit{om}}\nolimits}
\def\SheafExt{\mathop{\mathcal{E}\!\mathit{xt}}\nolimits}
\def\Sch{\mathit{Sch}}
\def\Mor{\mathop{\mathrm{Mor}}\nolimits}
\def\Ob{\mathop{\mathrm{Ob}}\nolimits}
\def\Sh{\mathop{\mathit{Sh}}\nolimits}
\def\NL{\mathop{N\!L}\nolimits}
\def\CH{\mathop{\mathrm{CH}}\nolimits}
\def\proetale{{pro\text{-}\acute{e}tale}}
\def\etale{{\acute{e}tale}}
\def\QCoh{\mathit{QCoh}}
\def\Ker{\mathop{\mathrm{Ker}}}
\def\Im{\mathop{\mathrm{Im}}}
\def\Coker{\mathop{\mathrm{Coker}}}
\def\Coim{\mathop{\mathrm{Coim}}}

% Boxtimes
%
\DeclareMathSymbol{\boxtimes}{\mathbin}{AMSa}{"02}

%
% Macros for moduli stacks/spaces
%
\def\QCohstack{\mathcal{QC}\!\mathit{oh}}
\def\Cohstack{\mathcal{C}\!\mathit{oh}}
\def\Spacesstack{\mathcal{S}\!\mathit{paces}}
\def\Quotfunctor{\mathrm{Quot}}
\def\Hilbfunctor{\mathrm{Hilb}}
\def\Curvesstack{\mathcal{C}\!\mathit{urves}}
\def\Polarizedstack{\mathcal{P}\!\mathit{olarized}}
\def\Complexesstack{\mathcal{C}\!\mathit{omplexes}}
% \Pic is the operator that assigns to X its picard group, usage \Pic(X)
% \Picardstack_{X/B} denotes the Picard stack of X over B
% \Picardfunctor_{X/B} denotes the Picard functor of X over B
\def\Pic{\mathop{\mathrm{Pic}}\nolimits}
\def\Picardstack{\mathcal{P}\!\mathit{ic}}
\def\Picardfunctor{\mathrm{Pic}}
\def\Deformationcategory{\mathcal{D}\!\mathit{ef}}

% --- appended 2026-10-12 (Chief): the project own definitions of the reference/slogan/history
% display environments, copied mechanically from the shared Stacks preamble lines 8-12
% (\usepackage{verbatim} plus the three \newenvironment lines, which the chapter preamble keeps
% verbatim from the project preamble). Nothing added, removed or redefined: these environments
% carry no mathematical content and the project itself typesets them as comments. Without these
% three lines a counted statement whose own text uses them fails with Environment slogan undefined.
\usepackage{verbatim}
\newenvironment{reference}{\comment}{\endcomment}
\newenvironment{slogan}{\comment}{\endcomment}
\newenvironment{history}{\comment}{\endcomment}
\usepackage{cleveref}
\crefname{originallemma}{original lemma}{original lemmas}
\Crefname{originallemma}{Original lemma}{Original lemmas}
\makeatletter
\newlabel{lemma-dominant-finite-number-irreducible-components}{{01RM}{source}{}{sourcepointer-1}{}}
\newlabel{lemma-dominant-finite-number-irreducible-components@cref}{{[originallemma][0][]01RM}{[1][0][]source}{}{}{}}
\makeatother
\crefname{originallemma}{original lemma}{original lemmas}
\Crefname{originallemma}{Original lemma}{Original lemmas}
\makeatletter
\newlabel{schemes-lemma-specialize-points}{{01J7}{source}{}{sourcepointer-2}{}}
\newlabel{schemes-lemma-specialize-points@cref}{{[originallemma][0][]01J7}{[1][0][]source}{}{}{}}
\makeatother
\title{Six equivalent conditions for dominance of a morphism between integral schemes}
\author{The Stacks Project Authors}
\date{stacks-project:morphisms.tex:local-snapshot-20261012; GNU FDL 1.2 or later}
\begin{document}
\maketitle
\noindent\textit{Editorial note.} All delivered mathematical text is the author's own, copied from the pinned Stacks Project chapter file pinned-original-source.tex; no statement or proof step was rewritten, reordered or removed.\par\medskip
\noindent\textit{Editorial note.} Following the multi-claim convention for this source, source\_url is the source-level URL https://stacks.math.columbia.edu/ and the source label lemma-dominant-between-integral is carried in claim\_key and here rather than in a per-claim tag URL.\par\medskip
\noindent\textit{Editorial note.} The preamble is the chapter's shared Stacks Project preamble, retained as sources/stacks-morphisms/retained-preamble.tex: a line-for-line excerpt of the shared Stacks preamble staged with this snapshot. No macro was invented or redefined, and this card adds no preamble text of its own.\par\medskip
\noindent\textit{Editorial note.} Statement span 1407-1425 and proof span 1427-1449 come from the halfA pool card operations/pool-cards/stacks-morphisms-halfA.json and were re-read in the pinned chapter file before this card was written.\par\medskip
\noindent\textit{Editorial note.} The proof of the second half says that one can argue 'in the exact same manner as above'; the argument for that half is nevertheless written out in the counted span (identification of local rings with spectra of points specializing to the point and the zero ideal reading), so the span is not a pointer.\par\medskip
\noindent\textit{Editorial note.} 2 cross-references occur inside the counted spans; they are listed in external\_references as navigation pointers only.\par\medskip
\noindent\textit{Allowed standard prerequisites (editorial).}\begin{itemize}
\item The same-chapter characterization of dominance by generic points of irreducible components, applied to an integral scheme whose irreducible component is the whole space. Cited at proof lines 1644-1645, this being the lane's card L4001 of this source. Gives the equivalence of dominance with mapping the generic point to the generic point.
\item The dictionary between points of the spectrum of a local ring and specializations of the point, cited from the schemes chapter. Cited at proof lines 1660-1662. Identifies the generic point of an integral scheme with the zero ideal of a local ring, which makes the local ring argument parallel to the affine chart argument.
\item Coordinate rings of nonempty affine opens in an integral scheme and local rings at points of it are integral domains. Used at proof lines 1646-1648 and 1654-1656. Lets injectivity of a ring map be tested by the preimage of the zero ideal.
\end{itemize}\medskip
\ifdefined\hypertarget\hypertarget{sourcepointer-1}{}\fi
\noindent\textit{Source reference.} Original lemma 01RM: focused-seeds/AGP/sources/stacks-snapshot-20261011/tags.txt, source lines 2307--2307. Same-chapter reference used in the retained author text at source line(s) 1429, 12858, 13423. It is supplied only as a navigation pointer so that the retained text compiles; the counted item delivers the author's own statement and the author's complete proof, and no part of the referenced item is copied in or substituted for the counted argument.\par\medskip
\ifdefined\hypertarget\hypertarget{sourcepointer-2}{}\fi
\noindent\textit{Source reference.} Original lemma 01J7: focused-seeds/AGP/sources/stacks-snapshot-20261011/tags.txt, source lines 2004--2004. Same-chapter reference used in the retained author text at source line(s) 1446, 12802, 13631, 13705, 13767. It is supplied only as a navigation pointer so that the retained text compiles; the counted item delivers the author's own statement and the author's complete proof, and no part of the referenced item is copied in or substituted for the counted argument.\par\medskip
\begin{lemma}
\label{lemma-dominant-between-integral}
Let $f : X \to Y$ be a morphism of integral schemes. The following
are equivalent
\begin{enumerate}
\item $f$ is dominant,
\item $f$ maps the generic point of $X$ to the generic point of $Y$,
\item for some nonempty affine opens $U \subset X$ and $V \subset Y$
with $f(U) \subset V$ the ring map $\mathcal{O}_Y(V) \to \mathcal{O}_X(U)$
is injective,
\item for all nonempty affine opens $U \subset X$ and $V \subset Y$
with $f(U) \subset V$ the ring map $\mathcal{O}_Y(V) \to \mathcal{O}_X(U)$
is injective,
\item for some $x \in X$ with image $y = f(x) \in Y$ the local ring
map $\mathcal{O}_{Y, y} \to \mathcal{O}_{X, x}$ is injective, and
\item for all $x \in X$ with image $y = f(x) \in Y$ the local ring
map $\mathcal{O}_{Y, y} \to \mathcal{O}_{X, x}$ is injective.
\end{enumerate}
\end{lemma}

\begin{proof}
The equivalence of (1) and (2) follows from
Lemma \ref{lemma-dominant-finite-number-irreducible-components}.
Let $U \subset X$ and $V \subset Y$ be nonempty affine opens with
$f(U) \subset V$. Recall that the rings $A = \mathcal{O}_X(U)$
and $B = \mathcal{O}_Y(V)$ are integral domains.
The morphism $f|_U : U \to V$ corresponds to a ring map
$\varphi : B \to A$. The generic points of $X$ and $Y$ correspond
to the prime ideals $(0) \subset A$ and $(0) \subset B$. Thus
(2) is equivalent to the condition $(0) = \varphi^{-1}((0))$,
i.e., to the condition that $\varphi$ is injective.
In this way we see that (2), (3), and (4) are equivalent.
Similarly, given $x$ and $y$ as in (5) the local rings
$\mathcal{O}_{X, x}$ and $\mathcal{O}_{Y, y}$ are domains
and the prime ideals $(0) \subset \mathcal{O}_{X, x}$
and $(0) \subset \mathcal{O}_{Y, y}$ correspond to the
generic points of $X$ and $Y$ (via the identification of
the spectrum of the local ring at $x$
with the set of points specializing to $x$, see
Schemes, Lemma \ref{schemes-lemma-specialize-points}).
Thus we can argue in the exact same manner as above
to see that (2), (5), and (6) are equivalent.
\end{proof}

\end{document}
